\documentclass[10pt,a4paper]{amsart}
\usepackage{amsmath,amssymb,amsthm,graphicx}
% Retained preamble for the Stacks Project chapter snapshot in this directory.
% The chapter's own preamble.tex is a copy of the shared Stacks Project preamble; the block
% below is sliced MECHANICALLY, line for line, from it. Nothing was added, removed or
% redefined; the wrapper class, the xr-hyper/hyperref wiring and the externaldocument list
% are omitted because the assembler supplies the document class.
%
% lines taken from preamble.tex: 3, 5, 9, 11, 12, 156, 157, 158, ... (51 lines total)
%
% package coverage restored from the chapter's own preamble (2026-10-13): xy, xy-2cell

\usepackage[all]{xy}\usepackage{verbatim}
\xyoption{2cell}
\UseAllTwocells

\newenvironment{reference}{\comment}{\endcomment}
\newenvironment{slogan}{\comment}{\endcomment}
\newenvironment{history}{\comment}{\endcomment}
\theoremstyle{plain}
\newtheorem{theorem}[subsection]{Theorem}
\newtheorem{proposition}[subsection]{Proposition}
\newtheorem{lemma}[subsection]{Lemma}
\theoremstyle{definition}
\newtheorem{definition}[subsection]{Definition}
\newtheorem{example}[subsection]{Example}
\newtheorem{exercise}[subsection]{Exercise}
\newtheorem{situation}[subsection]{Situation}
\theoremstyle{remark}
\newtheorem{remark}[subsection]{Remark}
\newtheorem{remarks}[subsection]{Remarks}
\numberwithin{equation}{subsection}
\def\lim{\mathop{\mathrm{lim}}\nolimits}
\def\colim{\mathop{\mathrm{colim}}\nolimits}
\def\Spec{\mathop{\mathrm{Spec}}}
\def\Hom{\mathop{\mathrm{Hom}}\nolimits}
\def\Ext{\mathop{\mathrm{Ext}}\nolimits}
\def\SheafHom{\mathop{\mathcal{H}\!\mathit{om}}\nolimits}
\def\SheafExt{\mathop{\mathcal{E}\!\mathit{xt}}\nolimits}
\def\Sch{\mathit{Sch}}
\def\Mor{\mathop{\mathrm{Mor}}\nolimits}
\def\Ob{\mathop{\mathrm{Ob}}\nolimits}
\def\Sh{\mathop{\mathit{Sh}}\nolimits}
\def\NL{\mathop{N\!L}\nolimits}
\def\CH{\mathop{\mathrm{CH}}\nolimits}
\def\proetale{{pro\text{-}\acute{e}tale}}
\def\etale{{\acute{e}tale}}
\def\QCoh{\mathit{QCoh}}
\def\Ker{\mathop{\mathrm{Ker}}}
\def\Im{\mathop{\mathrm{Im}}}
\def\Coker{\mathop{\mathrm{Coker}}}
\def\Coim{\mathop{\mathrm{Coim}}}
\DeclareMathSymbol{\boxtimes}{\mathbin}{AMSa}{"02}
\def\QCohstack{\mathcal{QC}\!\mathit{oh}}
\def\Cohstack{\mathcal{C}\!\mathit{oh}}
\def\Spacesstack{\mathcal{S}\!\mathit{paces}}
\def\Quotfunctor{\mathrm{Quot}}
\def\Hilbfunctor{\mathrm{Hilb}}
\def\Curvesstack{\mathcal{C}\!\mathit{urves}}
\def\Polarizedstack{\mathcal{P}\!\mathit{olarized}}
\def\Complexesstack{\mathcal{C}\!\mathit{omplexes}}
\def\Pic{\mathop{\mathrm{Pic}}\nolimits}
\def\Picardstack{\mathcal{P}\!\mathit{ic}}
\def\Picardfunctor{\mathrm{Pic}}
\def\Deformationcategory{\mathcal{D}\!\mathit{ef}}
\usepackage{cleveref}
\crefname{originallemma}{original lemma}{original lemmas}
\Crefname{originallemma}{Original lemma}{Original lemmas}
\makeatletter
\newlabel{perfect-lemma-kunneth}{{0FLQ}{source}{}{sourcepointer-1}{}}
\newlabel{perfect-lemma-kunneth@cref}{{[originallemma][0][]0FLQ}{[1][0][]source}{}{}{}}
\makeatother
\crefname{originallemma}{original lemma}{original lemmas}
\Crefname{originallemma}{Original lemma}{Original lemmas}
\makeatletter
\newlabel{sites-section-pullback}{{06UM}{source}{}{sourcepointer-2}{}}
\newlabel{sites-section-pullback@cref}{{[originallemma][0][]06UM}{[1][0][]source}{}{}{}}
\makeatother
\crefname{originallemma}{original lemma}{original lemmas}
\Crefname{originallemma}{Original lemma}{Original lemmas}
\makeatletter
\newlabel{stacks-cohomology-proposition-loc-qcoh-flat-base-change}{{0771}{source}{}{sourcepointer-3}{}}
\newlabel{stacks-cohomology-proposition-loc-qcoh-flat-base-change@cref}{{[originallemma][0][]0771}{[1][0][]source}{}{}{}}
\makeatother
\crefname{originallemma}{original lemma}{original lemmas}
\Crefname{originallemma}{Original lemma}{Original lemmas}
\makeatletter
\newlabel{stacks-cohomology-remark-QCoh-tensor}{{0GQM}{source}{}{sourcepointer-4}{}}
\newlabel{stacks-cohomology-remark-QCoh-tensor@cref}{{[originallemma][0][]0GQM}{[1][0][]source}{}{}{}}
\makeatother
\crefname{originallemma}{original lemma}{original lemmas}
\Crefname{originallemma}{Original lemma}{Original lemmas}
\makeatletter
\newlabel{stacks-cohomology-section-further-remarks}{{0GQS}{source}{}{sourcepointer-5}{}}
\newlabel{stacks-cohomology-section-further-remarks@cref}{{[originallemma][0][]0GQS}{[1][0][]source}{}{}{}}
\makeatother
\title{Pushforward product formula for flat representable morphisms}
\author{The Stacks Project Authors}
\date{stacks-project:stacks-more-morphisms.tex:local-snapshot-20261012; GNU FDL 1.2 or later}
\begin{document}
\maketitle
\noindent\textit{Editorial note.} Statement span is lines 4737-4747 of sources/stacks-stacks-more-morphisms/pinned-original-source.tex; proof span is lines 4749-4789, the complete author proof (39 body lines).\par\medskip
\noindent\textit{Editorial note.} Navigation stubs: 5 cited-but-not-included labels are supplied as non-load-bearing pointers with their pinned tags where the snapshot tag index has them.\par\medskip
\noindent\textit{Editorial note.} extra\_packages is empty on purpose: the retained preamble already loads xy with options and a second declaration clashes in the assembler wrapper.\par\medskip
\noindent\textit{Editorial note.} Text is retained verbatim, including the author's own typos and placeholders.\par\medskip
\noindent\textit{Editorial note.} Queue index 27 of queue\_stacks-stacks-more-morphisms.json (pb=1983, stmt 4737-4747, proof 4749-4789).\par\medskip
\noindent\textit{Editorial note.} stacks-stacks-more-morphisms wave smm01 (2026-10-14): fresh adjudication accept per lane-s3/smm\_judgments.md; pool per Chief 10:5x open-pool ruling (net 59 with-proof rows, bloodline lane-s3 by anchor audit, v2 true-key namespace stacks:stacks-more-morphisms.tex, row lemma-deform-property-fp-over-ft avoided on cross-chapter key collision); 60 with-proof rows judged 45 accept / 14 reject / 1 avoid; patch-sd-envs contract spans scanned.\par\medskip
\noindent\textit{Allowed standard prerequisites (editorial).}\begin{itemize}
\item Explicit prerequisite boundary. The cited inputs are prerequisites; the argument in the retained span is not replaced by any cited result.
\end{itemize}\medskip
\ifdefined\hypertarget\hypertarget{sourcepointer-1}{}\fi
\noindent\textit{Source reference.} Original lemma 0FLQ: sources/stacks-stacks-more-morphisms/pinned-original-source.tex, source lines 4789--4789. Cited by the retained statement or proof but not part of the collected spans; supplied only as a navigation pointer. The argument does not depend on its content.\par\medskip
\ifdefined\hypertarget\hypertarget{sourcepointer-2}{}\fi
\noindent\textit{Source reference.} Original lemma 06UM: sources/stacks-stacks-more-morphisms/pinned-original-source.tex, source lines 4789--4789. Cited by the retained statement or proof but not part of the collected spans; supplied only as a navigation pointer. The argument does not depend on its content.\par\medskip
\ifdefined\hypertarget\hypertarget{sourcepointer-3}{}\fi
\noindent\textit{Source reference.} Original lemma 0771: sources/stacks-stacks-more-morphisms/pinned-original-source.tex, source lines 4789--4789. Cited by the retained statement or proof but not part of the collected spans; supplied only as a navigation pointer. The argument does not depend on its content.\par\medskip
\ifdefined\hypertarget\hypertarget{sourcepointer-4}{}\fi
\noindent\textit{Source reference.} Original lemma 0GQM: sources/stacks-stacks-more-morphisms/pinned-original-source.tex, source lines 4789--4789. Cited by the retained statement or proof but not part of the collected spans; supplied only as a navigation pointer. The argument does not depend on its content.\par\medskip
\ifdefined\hypertarget\hypertarget{sourcepointer-5}{}\fi
\noindent\textit{Source reference.} Original lemma 0GQS: sources/stacks-stacks-more-morphisms/pinned-original-source.tex, source lines 4789--4789. Cited by the retained statement or proof but not part of the collected spans; supplied only as a navigation pointer. The argument does not depend on its content.\par\medskip
\begin{lemma}
\label{lemma-easy-kunneth}
Let $a : \mathcal{Y} \to \mathcal{X}$ and $b : \mathcal{Z} \to \mathcal{X}$
be representable by schemes, quasi-compact, quasi-separated, and flat.
Then $a_{\QCoh, *}\mathcal{O}_\mathcal{Y}
\otimes_{\mathcal{O}_\mathcal{X}}
b_{\QCoh, *}\mathcal{O}_\mathcal{Z} =
f_{\QCoh, *}\mathcal{O}_{\mathcal{Y} \times_\mathcal{X} \mathcal{Z}}$
where $f : \mathcal{Y} \times_\mathcal{X} \mathcal{Z} \to \mathcal{X}$
is the obvious morphism.
\end{lemma}

\begin{proof}
We abbreviate $\mathcal{P} = \mathcal{Y} \times_\mathcal{X} \mathcal{Z}$.
Since $a \circ \text{pr}_1 = f$ and $b \circ \text{pr}_2 = f$
we obtain maps $a_*\mathcal{O}_\mathcal{Y} \to f_*\mathcal{O}_\mathcal{P}$
and $b_*\mathcal{O}_\mathcal{Z} \to f_*\mathcal{O}_\mathcal{P}$ (using
relative pullback maps, see Sites, Section \ref{sites-section-pullback}).
Hence we obtain a relative cup product
$$
\mu :
a_*\mathcal{O}_\mathcal{Y}
\otimes_{\mathcal{O}_\mathcal{X}}
b_*\mathcal{O}_\mathcal{Z} \longrightarrow
f_*\mathcal{O}_{\mathcal{Y} \times_\mathcal{X} \mathcal{Z}}
$$
Applying $Q$ and its compatibility with tensor products
(Cohomology of Stacks, Remark \ref{stacks-cohomology-remark-QCoh-tensor})
we obtain an arrow
$Q(\mu) : a_{\QCoh, *}\mathcal{O}_\mathcal{Y}
\otimes_{\mathcal{O}_\mathcal{X}}
b_{\QCoh, *}\mathcal{O}_\mathcal{Z} \to
f_{\QCoh, *}\mathcal{O}_{\mathcal{Y} \times_\mathcal{X} \mathcal{Z}}$
in $\QCoh(\mathcal{O}_\mathcal{X})$.
Next, choose a scheme $U$ and a surjective smooth morphism
$U \to \mathcal{X}$. It suffices to prove the restriction of
$Q(\mu)$ to $U_\etale$ is an isomorphism, see
Cohomology of Stacks, Section \ref{stacks-cohomology-section-further-remarks}.
In turn, by the material in the same section,
it suffices to prove the restriction of $\mu$ to $U_\etale$
is an isomorphism (this uses that the source and target of $\mu$
are locally quasi-coherent modules with the base change property).
Moreover, we may compute pushforwards in the \'etale topology, see
Cohomology of Stacks, Proposition
\ref{stacks-cohomology-proposition-loc-qcoh-flat-base-change}.
Then finally, we see that
$a_*\mathcal{O}_\mathcal{Y}|_{U_\etale} = (V \to U)_{small, *}\mathcal{O}_V$
where $V = U \times_\mathcal{X} \mathcal{Y}$.
Similarly for $b_*$ and $f_*$. Thus the result follows
from the K\"unneth formula for flat, quasi-compact, quasi-separated
morphisms of schemes, see Derived Categories of Schemes,
Lemma \ref{perfect-lemma-kunneth}.
\end{proof}

\end{document}
